\section{Operational implications}
\label{sec:implications}

Each statement below restates a measured result; the controllers of
Section~\ref{sec:mitigation} are model-tested only, no controller has been validated on
hardware, and no numeric threshold transfers beyond the measured
stack.

\begin{itemize}
\item Horizon rule. Benchmarks shorter than $\sim$300\,min overstate stability
  at boundary operating points: 2 of 8 surviving runs at
  ($N{=}8$, 0.020) collapsed by the delayed mechanism near minute 240,
  past every 60--180\,min window (Section~\ref{sec:horizon}) \evid{R3}.
\item Routing affinity is a stability parameter, not a latency
  parameter. In the one measured run ($n=1$), the load-only
  configuration ran warm at $h$ 0.53 versus 0.94 and collapsed after
  the surge with no recovery phase (Section~\ref{sec:affinity}) \evid{R14}.
\item Router span is a blast-radius parameter. Every measured 16-replica
  collapse was fleet-total; independent 8-replica shard routers
  contained every collapse inside the failing partition
  (Section~\ref{sec:shards}) \evid{R10}. Sharding is a measured containment boundary.
\item Drain-depth observables are early-warning candidates. $D$ separates
  immediate collapse from other outcomes in all 8 runs at the
  boundary point and adds signal beyond rate over 21 runs
  (Section~\ref{sec:separatrix}) \evid{R15}. $D$ is a mediating observable whose threshold
  moves with rate and span and which does not predict the delayed
  class; unlike a fixed-constant admission threshold
  \cite{chen2026concur}, deployment use would require local (rate, $N$)
  calibration.
\item Tier sizing is a persistence decision, not only a hit-rate
  decision. At ($N{=}8$, 0.022) the smallest tier falls through
  congestion to cold collapse ($n=1$) while the two larger sizes never
  go cold in 8 runs (Section~\ref{sec:sizeseries}); capacity sets regime persistence
  per rule~\eqref{eq:ccpustar}, and the pinned restore channel sets congested-state
  throughput per bound~\eqref{eq:murestore}, with restore bandwidth itself unvaried
  \evid{R11}.
\end{itemize}

One model-supported implication (Section~\ref{sec:mitigation}, not measured): admission
control for this failure should key on basin coordinates - KV
occupancy and backlog depth - and defer rather than reject.
Controllers keyed on the collapse symptoms themselves (hit rate,
queueing delay) engaged only after the basin boundary was crossed in
the model screen, and rejection spent permanent losses at operating
points that mostly recover unaided.
